\documentclass[aspectratio=169,11pt]{beamer}

\usetheme{Madrid}
\usecolortheme{default}
\usefonttheme{professionalfonts}
\setbeamertemplate{navigation symbols}{}
\setbeamertemplate{footline}[frame number]

\usepackage{amsmath, amssymb, booktabs}

\newcommand{\E}{\mathbb{E}}
\newcommand{\1}{\mathbf{1}}
\newcommand{\Prb}{\mathbb{P}}

\title{Causal Inference With Network Data}
\subtitle{Empirical Methods --- Week 19}
\author{Teaching deck draft}
\date{}

\begin{document}

\begin{frame}
  \titlepage
\end{frame}

\begin{frame}{Week 19 question}
\begin{itemize}
    \item How can researchers estimate effects when units influence one another through observed or latent links?
    \item The problem is causal inference, not general network analysis.
    \item Network structure changes identification, estimation, inference, and interpretation.
    \item The graph definition is substantive: it states who can affect whom and through which channel.
\end{itemize}

\medskip
Lecture 18 built the graph. Lecture 19 asks what can be learned causally once the graph exists.
\end{frame}

\begin{frame}{Why Lectures 1--5 are not always enough}
\small
\begin{tabular}{p{0.28\linewidth} p{0.30\linewidth} p{0.30\linewidth}}
\toprule
Core design habit & Network complication & Response \\
\midrule
Individual treatment contrast & untreated units may be exposed & define exposure states \\
SUTVA notation & outcomes depend on others' treatment & potential outcomes under interference \\
Peer regression & outcomes jointly determine each other & topology, timing, randomization \\
One row per unit & rows are dyads or links & dyadic/network inference \\
Cluster standard errors & dependence follows shared nodes & pair-aware variance logic \\
\bottomrule
\end{tabular}

\medskip
The core toolkit still matters; the network tells us when its assumptions are too small.
\end{frame}

\begin{frame}{Potential outcomes under interference}
\small
Standard binary-treatment notation:
\[
\tau=\E[Y_i(1)-Y_i(0)]
\]

Network notation:
\[
Y_i = Y_i(Z_i,Z_{-i};G)
\]

Exposure-mapping restriction:
\[
D_i=g_i(Z_{-i},G), \qquad Y_i=Y_i(Z_i,D_i)
\]

\begin{itemize}
    \item \(G\): graph of economically relevant links.
    \item \(Z_i\): own treatment status.
    \item \(D_i\): exposure to others' treatments through the graph.
    \item The mapping \(g_i(\cdot)\) is an identifying assumption, not a formatting choice.
\end{itemize}
\end{frame}

\begin{frame}{The reflection problem}
\small
\[
Y_i
=
\alpha
+ \beta \sum_j w_{ij}Y_j
+ \gamma \sum_j w_{ij}X_j
+ \delta X_i
+ \varepsilon_i
\]

\begin{itemize}
    \item \(\beta\): endogenous peer effect.
    \item \(\gamma\): contextual peer-characteristic effect.
    \item Correlated effects: sorting, common shocks, shared managers, shared neighborhoods.
    \item Peer outcomes and own outcomes can be jointly determined.
    \item Leave-one-out peer means reduce mechanical self-inclusion; they do not solve sorting.
\end{itemize}

\medskip
Network topology helps only when the exclusion logic is economically credible.
\end{frame}

\begin{frame}{Interference and exposure mappings}
\small
\[
\tau^{direct}(d)=\E[Y_i(1,d)-Y_i(0,d)]
\]
\[
\tau^{spill}(z;d,d')=\E[Y_i(z,d)-Y_i(z,d')]
\]

\begin{itemize}
    \item Direct effects compare own treatment at a fixed exposure state.
    \item Spillover effects compare exposure states at a fixed own treatment.
    \item Exposure can be treated-neighbor share, treated manager, treated high-centrality contact, or saturation cell.
    \item Failure mode: the exposure mapping misses the true path or compresses heterogeneous paths too aggressively.
\end{itemize}
\end{frame}

\begin{frame}{Randomized saturation}
\small
\[
Z_i \sim Bernoulli(p_c), \qquad p_c \in \{p_L,p_H\}
\]

\begin{itemize}
    \item Randomize treatment intensity across teams, classrooms, villages, worksites, or graph partitions.
    \item Untreated units in high-saturation clusters reveal spillovers relative to untreated units in low-saturation clusters.
    \item Treated units across saturation cells reveal whether direct effects depend on peer exposure.
    \item Partial interference assumes spillovers stay inside the cluster.
    \item Strong design, but fragile if spillovers cross the boundary.
\end{itemize}
\end{frame}

\begin{frame}{Network-targeted econometric methods}
\small
\begin{tabular}{p{0.28\linewidth} p{0.31\linewidth} p{0.29\linewidth}}
\toprule
Tool & Identifies or protects & Main caveat \\
\midrule
Topology-based peer design & peer effects under excluded-peer variation & network and exclusion restrictions \\
Exposure mapping & effects of direct and indirect exposure states & mapping must match mechanism \\
Saturation experiment & direct, spillover, intensity effects & boundaries must be plausible \\
Randomization inference & design-based p-values & must follow assignment rule \\
Leave-one-out peers & cleaner peer exposure measure & not a full design \\
Dyadic inference & valid pair-level uncertainty & does not solve link endogeneity \\
\bottomrule
\end{tabular}
\end{frame}

\begin{frame}{Dyadic inference}
\small
Dyadic model:
\[
Y_{ij}=X_{ij}'\beta+a_i+b_j+u_{ij}
\]

\begin{itemize}
    \item Observations are pairs: worker--firm, employee--applicant, worker--worker, sender--receiver.
    \item Dyads sharing \(i\) can be correlated.
    \item Dyads sharing \(j\) can be correlated.
    \item IID or ordinary robust standard errors usually understate uncertainty.
    \item Pair-aware inference clusters or otherwise accounts for shared-node dependence.
\end{itemize}

\medskip
The estimating unit and variance formula must match the relational data structure.
\end{frame}

\begin{frame}{Network definition changes interpretation}
\small
\begin{tabular}{p{0.24\linewidth} p{0.34\linewidth} p{0.30\linewidth}}
\toprule
Link & Possible mechanism & Interpretation risk \\
\midrule
Coworker & learning, norms, monitoring, team production & same firm may just be sorting \\
Manager & supervision, information, practices & manager assignment may be endogenous \\
Referral & information, screening, trust, favoritism & tie may proxy applicant quality \\
Neighborhood & job information, local norms, commute knowledge & proximity may proxy local shocks \\
Platform interaction & congestion, reputation, matching & platform rules shape observed links \\
\bottomrule
\end{tabular}

\medskip
The graph is part of the estimand, not just an input file.
\end{frame}

\begin{frame}{Paper anchors}
\begin{itemize}
    \item Manski: reflection problem and endogenous social effects.
    \item Bramoulle, Djebbari, Fortin: network topology and peer-effect identification.
    \item Aronow and Samii: causal effects under general interference.
    \item Athey, Eckles, Imbens: exact inference with network interference.
    \item Cornelissen, Dustmann, Schonberg: workplace peer effects.
    \item Pallais and Sands: referrals and field-experimental labor-market design.
    \item Graham and Canen: dyadic/network data and dependence-aware inference.
\end{itemize}
\end{frame}

\begin{frame}{Research Lab design}
\begin{columns}[T,onlytextwidth]
\begin{column}{0.32\linewidth}
\textbf{Reproduce}
\begin{itemize}
    \item synthetic teams
    \item coworker links
    \item treatment saturation
    \item leave-one-out peers
    \item spillover estimates
\end{itemize}
\end{column}
\begin{column}{0.32\linewidth}
\textbf{Diagnose}
\begin{itemize}
    \item exposure mappings
    \item reflection risk
    \item randomization inference
    \item partial interference
    \item common shocks
\end{itemize}
\end{column}
\begin{column}{0.32\linewidth}
\textbf{Transfer}
\begin{itemize}
    \item referral dyads
    \item pair-level model
    \item naive vs dyadic SE
    \item access inequality
    \item design memo
\end{itemize}
\end{column}
\end{columns}
\end{frame}

\begin{frame}{What each method identifies}
\small
\begin{tabular}{p{0.25\linewidth} p{0.34\linewidth} p{0.29\linewidth}}
\toprule
Method & Target & Not automatically identified \\
\midrule
Peer-effect regression & influence within a defined group & formation or common shocks \\
Exposure mapping & effect of exposure state \(d\) & all paths in the network \\
Saturation design & spillovers within design cells & market-wide equilibrium effects \\
Network experiment & effect under assignment protocol & effects under a different graph \\
Dyadic model & pair-level association or effect & why the link exists \\
\bottomrule
\end{tabular}
\end{frame}

\begin{frame}{Takeaways}
\begin{itemize}
    \item Network causal inference starts by defining who can affect whom.
    \item Reflection separates peer correlations from peer effects.
    \item Exposure mappings turn impossible potential-outcome objects into estimands.
    \item Saturation and graph experiments randomize exposure, not only own treatment.
    \item Partial interference is an assumption about boundaries.
    \item Dyadic data need dependence-aware inference.
    \item Lecture 20 asks what happens when the network itself is endogenous.
\end{itemize}
\end{frame}

\end{document}
